% !TeX program = tectonic
\documentclass[11pt,a4paper]{article}
\usepackage{fontspec}
\usepackage[russian]{babel}
\babelfont{rm}[Language=Default]{Liberation Serif}
\babelfont[Language=Default]{sf}{Carlito}
\babelfont[Language=Default]{tt}{DejaVu Sans Mono}
\usepackage{amsmath,amssymb,amsthm}
\usepackage{geometry}
\geometry{a4paper, top=2.4cm, bottom=2.4cm, left=2.4cm, right=2.4cm}
\usepackage{booktabs,tabularx,enumitem,xcolor,tcolorbox}
\tcbuselibrary{breakable,skins}
\definecolor{linkblue}{HTML}{1F4E79}
\definecolor{statgreen}{HTML}{2F6B3A}
\definecolor{goldacc}{HTML}{8B7E5A}
\usepackage[colorlinks=true,linkcolor=linkblue,urlcolor=linkblue,unicode=true]{hyperref}
\theoremstyle{plain}
\newtheorem{theorem}{Теорема}
\newtheorem{lemma}[theorem]{Лемма}
\newtheorem{proposition}[theorem]{Предложение}
\newtheorem{corollary}[theorem]{Следствие}
\theoremstyle{definition}
\newtheorem{definition}[theorem]{Определение}
\theoremstyle{remark}
\newtheorem{remark}[theorem]{Замечание}
\newtcolorbox{statusbox}[1][]{colback=statgreen!5,colframe=statgreen!75,
  title={\textbf{Сводка доказанного:} #1},fonttitle=\small,boxrule=0.7pt,arc=2pt,breakable}
\newtcolorbox{databox}[1][]{colback=goldacc!6,colframe=goldacc!80,
  title={\textbf{Данные протокола:} #1},fonttitle=\small,boxrule=0.7pt,arc=2pt,breakable}
\newtcolorbox{verifybox}[1][]{colback=linkblue!5,colframe=linkblue!80,
  title={\textbf{Верификация:} #1},fonttitle=\small,boxrule=0.7pt,arc=2pt,breakable}
\widowpenalty=10000
\clubpenalty=10000
\setlength{\parskip}{0.35em}
\newcommand{\CC}{\mathbb{C}}
\newcommand{\RR}{\mathbb{R}}
\newcommand{\ZZ}{\mathbb{Z}}
\newcommand{\QQ}{\mathbb{Q}}
\newcommand{\Ga}{\Gamma}
\newcommand{\Bfun}{\mathrm{B}}
\newcommand{\im}{\operatorname{Im}}
\newcommand{\re}{\operatorname{Re}}
\newcommand{\lcm}{\operatorname{lcm}}
\newcommand{\SNF}{\operatorname{SNF}}
\newcommand{\Jac}{\operatorname{Jac}}
\newcommand{\rank}{\operatorname{rank}}
\newcommand{\Ind}{\operatorname{Index}}
\newcommand{\disc}{\operatorname{disc}}
\begin{document}
\begin{center}
{\LARGE\bfseries Теорема 3: поляризация и индексы решётки Ходжа}\\[0.4em]
{\large главная поляризация и вычислимость индексов по блокам}\\[0.6em]
{\normalsize Исаев Исхак Хамзатович}\\[0.2em]
{\small Программа <<Динамический принцип>> --- лаборатория \texttt{hodge-laboratory}}\\[0.15em]
{\small Отдельная монография теоремы · Монография 6/16}
\end{center}
\vspace{0.4em}
{\color{linkblue}\hrule height 0.8pt}
\vspace{0.8em}

\section{Постановка}

Форма пересечений $J$ на $H_1(C_N,\ZZ)$ --- не просто топологический
инвариант: она является поляризацией структуры Ходжа веса $1$ и
позволяет определить индексы решётки Ходжа по блокам проводников ---
целочисленные инварианты, вычислимые алгоритмически.

\begin{definition}[поляризация]
Для замкнутых форм $\widetilde E([\alpha],[\beta])
=\iint_{C_N}\alpha\wedge\beta$; поляризационная форма на
$H_1(C_N,\ZZ)$ --- форма пересечений $E:=J$.
\end{definition}

\begin{theorem}[$J$ --- поляризация веса 1]\label{th:pol}
Относительно разложения Ходжа
$H^1=H^{1,0}\oplus H^{0,1}$ выполнены соотношения Римана--Ходжа:
(1) $\widetilde E(\omega,\omega')=0$ для голоморфных форм;
(2) $-\frac1{2i}\widetilde E(\omega,\overline{\omega})
=\int_{C_N}|g|^2dxdy>0$ для $0\ne\omega=g\,dz$.
Следовательно, $\Jac(C_N)=H^1/H_1$ с поляризацией $J$ ---
главное поляризованное абелево многообразие.
\end{theorem}

\begin{proof}
(i) На кривой нет ненулевых голоморфных 2-форм:
$\Omega^2_{C_N}=0$, поэтому $\omega\wedge\omega'=0$ --- это первое
соотношение Римана, прочитанное на билинейном тождестве.
(ii) $\omega\wedge\overline{\omega}=-2i|g|^2dx\wedge dy$, откуда
$-\frac1{2i}\widetilde E(\omega,\overline{\omega})
=\int_{C_N}|g|^2dxdy>0$ при $\omega\ne0$.
Унимодулярность $J$ (симплектический базис) даёт главность
поляризации.
\end{proof}

\begin{definition}[индексы по блокам]
По проводнику $d\mid N$: решётка блока
$\Lambda_d$ --- изотипная компонента характеров с проводником $d$;
насыщение $\Lambda_d^{sat}$; индекс блока
$\Ind(\Lambda_d)=[\Lambda_d^{sat}:\Lambda_d]$. Данные блока: ранг,
сигнатура $E|_{\Lambda_d}$, определитель, индекс.
\end{definition}

\begin{theorem}[вычислимость индексов]\label{th:indcomp}
Все величины конечны и вычисляются алгоритмически из замкнутых форм:
насыщение --- приведением Эрмита--Смита; сигнатура и определитель ---
из комбинаторной матрицы пересечений с фильтрацией по характерам;
индекс --- отношение определителей. Число собственных форм по блокам
даёт перепись $h_d$: $1+6+84=91$ ($N=15$);
$1+6+9+30+84+276=406$ ($N=30$).
\end{theorem}

\begin{proof}
Конечность: $H_1(C_N,\ZZ)$ --- свободная абелева группа ранга $2g$.
Вычислимость: изотипная компонента определена над $\QQ(\zeta_N)$;
её пересечение с целочисленной решёткой --- целочисленная
подрешётка; базис насыщения --- алгоритм Эрмита--Смита по таблице
действия $\mu_N\times\mu_N$ (комбинаторика листов). Ограничение
поляризации: $E=J$, матрица комбинаторна. Перепись $h_d$ ---
приложение монографии; суммы совпадают с родом (V1).
\end{proof}

\begin{databox}[числа --- сертификат H]
Тип поляризации стенда (SNF): $(1,1,5,5,15,15,15,15)$;
произведение инвариантных множителей $=1265625=3^4\cdot5^6
=1125^2$ --- целочисленное согласование с дискриминантом;
$\mathrm{vol}_h=1125$.
\end{databox}

\begin{statusbox}[итог]
Поляризация доказана (Римана--Ходжа); вычислимость индексов
доказана алгоритмически; сертификат H реализует вычисление с точной
целочисленной сертификацией --- индексы не только вычислимы, но и
вычислены.
\end{statusbox}

\begin{verifybox}[как воспроизвести]
\texttt{python3 laboratory.py} $\to$ <<Сертификаты $\to$ H1>>;
Lean: \texttt{snf\_product} и \texttt{stand\_disc} --- машинные
доказательства $1265625=1125^2$ и произведения SNF.
\end{verifybox}

\vfill
{\color{linkblue}\hrule height 0.6pt}
\vspace{0.5em}
{\small\color{gray}
Лицензия: индивидуальная исключительная лицензия Исаева Исхака Хамзатовича.
Все права на вычисления, формулы и тексты принадлежат автору.
Верификация: \texttt{python3 laboratory.py} (пункт <<Стенды>>/<<Сертификаты>>),
Lean: \texttt{verification/lean/HodgeLaboratory.lean}.
}
\end{document}
